We asked whether fresh, firm-level earnings-call sentiment improves cross-sectional prediction of monthly excess stock returns. Within our framework, the answer is a qualified yes. The preprocessing ablation (Table~\ref{tab:preproc}) shows that freshness filtering, delta features, and cross-sectional z-scoring are each necessary: freshness filtering provides the first profitable baseline and z-scoring lifts Sharpe from 0.14 to 0.79. Over the 2019--2023 test period, FinBERT-only XGBoost achieves $+11.4\%$ net annualised return (Sharpe 0.79) and a Rank IC $t$-statistic of 3.79 over the 108-month post-training window. A cross-sectional MLP on the same features produces $+7.87\%$ net (Sharpe 0.58), and a Ridge baseline (Sharpe $-0.10$) suggests the relationship is not well captured by a simple linear specification.

JKP factor returns are not well aligned with this cross-sectional strategy design: being constant across stocks within a month, they encode market-level time-series information rather than firm-level cross-sectional differences. Sequence models (LSTM, Transformer) are not robust out-of-sample because quarterly earnings-call trajectories are sparse and regime-dependent.

We caution that the signal is modest (monthly IC $\approx 0.03$), the test period is limited (60 months), and the covered universe is biased toward mid- and large-cap firms with regular earnings calls. These results are consistent with the hypothesis that fresh changes in earnings-call tone contain incremental information not fully captured by our benchmarks. Future work could explore higher-frequency signal updates, alternative language models, or broader universe coverage to assess generalisability.
